% biblatex-spbasic documentation 2021/07/11 
%
% Copyright (c) 2021 Herbert Voß hvoss@tug.org
%% $Id: biblatex-spbasic.tex 379 2021-12-29 13:38:14Z herbert $

\documentclass[toc=graduated,parskip=half-,DIV=13,fontsize=11pt,paper=a4]{scrartcl}

\usepackage{eurosym}	% Damit das Euro-Symbol dargestellt werden kann
\usepackage{libertinus-otf} 		% Damit LaTex Umlaute usw. erkennt
\setmonofont[Scale=MatchLowercase,FakeStretch=0.9]{AnonymousPro}
\usepackage[english]{babel}
\usepackage{listings}
\lstset{basicstyle=\ttfamily}
\usepackage{hvlogos}
\usepackage{xcolor}
\usepackage{enumitem}
\setlist{nosep}

\usepackage[babel]{csquotes}
\usepackage{xurl}
\usepackage[style=biblatex-spbasic,hyperref]{biblatex} 

\usepackage{xspace}
\addbibresource{biblatex-examples.bib}
\addbibresource{\jobname.bib}

\def\bibsp{\texttt{biblatex-spbasic}\xspace}
\def\bib{\texttt{biblatex}\xspace}
\def\cs#1{\texttt{\textbackslash#1}}

\usepackage[colorlinks,linktocpage,citecolor=blue!100!black!60]{hyperref}

\begin{document}

\title{biblatex-spbasic}
\subtitle{Ver. 0.04}
\author{Herbert Voß}

\maketitle

\vfill
Thanks to: Wolfgang Engelmann; Moritz Wemheuer;
\newpage

\section{Introduction}
This documentation briefly explains the adaptation of the \texttt{biblatex-spbasic}
bibliography style and how it can be used. Some of the functions of \texttt{biblatex} 
that are based on are also indirectly explained. If you have any questions about the 
functions of \texttt{biblatex}, you should always consult the very good and detailed 
\texttt{biblatex} documentation. 

%Biblatex and biblatex-spbasic cannot be used meaningfully without at least skimming through the user guide in the biblatex documentation.
For a quick start, the requirements for using \texttt{biblatex-spbasic} are briefly explained in the following section.


\section{Installation}

If you do have an official installation of \MiKTeX\ or \TeXLive, then everything would be there anyway, otherwise
use the package manager and install \texttt{biblatex-spbasic}. If you do not have such an official distribution
installed, then copy the three files

\begin{verbatim}
biblatex-spbasic.bbx
biblatex-spbasic.cbx
biblatex-spbasic.lbx
\end{verbatim}

into your local \TeX\ directory and update the so-called file name database by running \texttt{texhash}.

\section{Using the style}

\texttt{biblatex-spbasic} is based on the packages from \texttt{biblatex} and especially the style \texttt{authoryear}. 
The reason for this is that \texttt{biblatex} provides all functions that make it possible to develop individual 
styles. However, the way in which »Springer Verlag« cites is so special compared to what is otherwise usual 
that additional functions are required to map them completely. These can be found in \texttt{authoryear.bbx}. 
You could have taken these out of the package and integrated them into \texttt{biblatex-spbasic}, but the 
effort involved seemed too great to me. In addition, the well-maintained \texttt{biblatex} package ensures 
that the functions provided by it are reliably retained and compatible despite the very rapid development of \texttt{biblatex}.

It is obvious that using \texttt{biblatex-spbasic} needs the \texttt{biblatex} package, too. 
\texttt{biblatex-spbasic} is loaded by

\begin{verbatim}
\usepackage[style=biblatex-spbasic,pubstate=<true|false>]{biblatex}
\end{verbatim}

It is possible to load additional optional arguments which refer to \texttt{biblatex}, the package
\texttt{biblatex-spbasic} itself has only the option \texttt{pubstate}, which is set by default to false. There
is no need to set it, if you want this presetting. For printing the publication state use \texttt{pubstate=true}.

\section{Example}
This is some senseless\footfullcite{Young2018}  text but with a citation (\verb|\parencite|).~\parencite[See][pp.~16]{Daan1984}
And this is a \verb|\textcite|, what \textcite[See][pp.~16]{Young2018} already wrote. Remember, this
text is nonsense for this example!\footcite{Young2018}  And the 
authors~\textcite{Darwin1880} wrote an article titled~\citetitle{Darwin1880}.

\nocite{*}
\printbibliography[title=Bibliography]


\end{document}

